\section*{Original introductory normalizations}
This paper deals with maximal representations of complex hyperbolic lattices in semisimple Hermitian Lie groups with no compact factors.

A complex hyperbolic lattice $\Gamma$ is a lattice in the Lie group $\mathrm{SU}(1,n)$, a finite cover of the group of biholomorphisms of the $n$-dimensional complex hyperbolic space $\mathbb H^n_{\mathbb C}=\mathrm{SU}(1,n)/\mathrm U(n)$. We shall always assume that our lattice $\Gamma$ is cocompact, or uniform, meaning that the quotient $X:=\Gamma\backslash\mathbb H^n_{\mathbb C}$ is compact. To simplify matters, we also assume in this introduction (except in the statements of our main results) that $\Gamma$ is torsion free, so that $X$ is also a manifold. The $\mathrm{SU}(1,n)$-invariant Kähler form on $\mathbb H^n_{\mathbb C}$ with constant holomorphic sectional curvature $-1$ and the Kähler form it induces on $X$ will both be denoted by $\omega$.

A Lie group $G_{\mathbb R}$ is said to be a \emph{real algebraic Hermitian Lie group} (or a Hermitian group for short) if it is the connected component $\jepPubBmi{G}(\mathbb R)^{\circ}$ of the group of real point of an algebraic group $\jepPubBmi{G}$ defined over $\mathbb R$, if it is semisimple with no compact factors, and if its associated symmetric space $M=G_{\mathbb R}/K_{\mathbb R}$ is a Hermitian symmetric space (of the noncompact type). This means that the noncompact symmetric space $M$ admits a $G_{\mathbb R}$-invariant complex structure, with respect to which the $G_{\mathbb R}$-invariant Riemannian metrics are (necessarily) Kähler. The real rank $\operatorname{rk}_{\mathbb R}G_{\mathbb R}$ of $G_{\mathbb R}$ coincides with the rank $r_M$ of $M$ as a symmetric space, namely the maximal dimension of a flat subspace in $M$. Simple Hermitian groups are classified: there are four infinite families of classical groups, which up to isogeny are $\mathrm{SU}(p,q)$ with $1\leqslant p\leqslant q$, $\mathrm{SO}_0(2,p)$ with $p\geqslant3$, $\mathrm{Sp}(2m,\mathbb R)$ with $m\geqslant2$ and $\mathrm{SO}^{\star}(2m)$ with $m\geqslant4$; and two exceptional groups $\mathrm E_{6(-14)}$ and $\mathrm E_{7(-25)}$. The real ranks of these groups are respectively $p$, $2$, $m$, $\lfloor m/2\rfloor$, $2$ and $3$.

Let $\omega_M$ denote the $G_{\mathbb R}$-invariant Kähler form of $M$, uniquely normalized so that its holomorphic sectional curvatures lie between $-1$ and $-1/r_M$.

If $\rho$ is a representation (a group homomorphism) from $\Gamma$ to $G_{\mathbb R}$, we define its \emph{Toledo invariant} $\tau(\rho)$ as follows:
\[
\tau(\rho)=\frac1{n!}\int_X f^{\star}\omega_M\wedge\omega^{n-1},
\]
where $f:\mathbb H^n_{\mathbb C}\to M$ is a $\jepPubScript{C}^{\infty}$ and $\rho$-equivariant map and $f^{\star}\omega_M$ is seen as a 2-form on $X$ by $\Gamma$-invariance. The Toledo invariant does not depend on the choice of the map $f$, it depends only on $\rho$, and in fact, only on the connected component of $\rho$ in $\operatorname{Hom}(\Gamma,G_{\mathbb R})$. Moreover, it satisfies the following Milnor-Wood type inequality:
\[
|\tau(\rho)|\leqslant r_M\operatorname{vol}(X),
\]
a fundamental property established in full generality in \cite{jep93BI07}.
The \emph{maximal representations} $\rho:\Gamma\to G_{\mathbb R}$ are those representations for which the Milnor-Wood inequality is an equality.

One interesting by-product of this unified approach is that it allows to exclude a priori the possibility of maximal representations in any tube type real algebraic Hermitian Lie group, and in particular in $\operatorname E_{7(-25)}$. Recall that up to isogeny the simple tube type Hermitian groups are $\operatorname{SU}(p,p)$, $\operatorname{SO}_0(2,p)$, $\operatorname{Sp}(2m,\mathbb R)$, $\operatorname{SO}^{\star}(2m)$ with $m$ even, and $\operatorname E_{7(-25)}$. See also Section~2.4.4. Indeed, we prove in Section~5.1 that for tube type targets the representation $\rho$ satisfies a stronger inequality than the Milnor-Wood inequality.
\begin{propositionC}
Let $\Gamma$ be a uniform lattice in $\mathrm{SU}(1,n)$, with $n\geqslant2$, and let $X=\Gamma\backslash\mathbb H^n_{\mathbb C}$. Assume that the real algebraic Hermitian Lie group $G_{\mathbb R}$ has tube type and let $r_M$ be the real rank of $G_{\mathbb R}$. Let $\rho$ be a representation $\Gamma\to G_{\mathbb R}$. Then
\[
|\tau(\rho)|\leqslant\max\Big\{r_M-1,\frac{r_M}{2}\cdot\frac{n+1}{n}\Big\}\operatorname{vol}(X)<r_M\operatorname{vol}(X).
\]
\end{propositionC}

